% !Mode:: "TeX:UTF-8"
\documentclass[12pt,a4paper]{article}

%%%%%%%%------------------------------------------------------------------------
%%%% 日常所用宏包

%% 控制页边距
% 如果是beamer文档类, 则不用geometry
\makeatletter
\@ifclassloaded{beamer}{}{\usepackage[top=2.5cm, bottom=2.5cm, left=2.5cm, right=2.5cm]{geometry}}
\makeatother

%% 控制项目列表
\usepackage{enumerate}

%% 多栏显示
\usepackage{multicol}

%% 算法环境
\usepackage{algorithm}  
\usepackage{algorithmic} 
\usepackage{float} 

%% 网址引用
\usepackage{url}

%% 控制矩阵行距
\renewcommand\arraystretch{1.4}

%% hyperref宏包，生成可定位点击的超链接，并且会生成pdf书签
\makeatletter
\@ifclassloaded{beamer}{
\usepackage{hyperref}
\usepackage{ragged2e} % 对齐
}{
\usepackage[%
    pdfstartview=FitH,%
    CJKbookmarks=true,%
    bookmarks=true,%
    bookmarksnumbered=true,%
    bookmarksopen=true,%
    colorlinks=true,%
    citecolor=blue,%
    linkcolor=blue,%
    anchorcolor=green,%
    urlcolor=blue%
]{hyperref}
}
\makeatother



\makeatletter % 如果是 beamer 不需要下面两个包
\@ifclassloaded{beamer}{
\mode<presentation>
{
} 
}{
%% 控制标题
\usepackage{titlesec}
%% 控制目录
\usepackage{titletoc}
}
\makeatother

%% 控制表格样式
\usepackage{booktabs}

%% 控制字体大小
\usepackage{type1cm}

%% 首行缩进，用\noindent取消某段缩进
\usepackage{indentfirst}

%% 支持彩色文本、底色、文本框等
\usepackage{color,xcolor}

%% AMS LaTeX宏包: http://zzg34b.w3.c361.com/package/maths.htm#amssymb
\usepackage{amsmath,amssymb}
%% 多个图形并排
\usepackage{subfig}
%%%% 基本插图方法
%% 图形宏包
\usepackage{graphicx}
\newcommand{\red}[1]{\textcolor{red}{#1}}
\newcommand{\blue}[1]{\structure{#1}}
\newcommand{\brown}[1]{\textcolor{brown}{#1}}
\newcommand{\green}[1]{\textcolor{green}{#1}}


%%%% 基本插图方法结束

%%%% pgf/tikz绘图宏包设置
\usepackage{pgf,tikz}
\usetikzlibrary{shapes,automata,snakes,backgrounds,arrows}
\usetikzlibrary{mindmap}
%% 可以直接在latex文档中使用graphviz/dot语言，
%% 也可以用dot2tex工具将dot文件转换成tex文件再include进来
%% \usepackage[shell,pgf,outputdir={docgraphs/}]{dot2texi}
%%%% pgf/tikz设置结束


\makeatletter % 如果是 beamer 不需要下面两个包
\@ifclassloaded{beamer}{

}{
%%%% fancyhdr设置页眉页脚
%% 页眉页脚宏包
\usepackage{fancyhdr}
%% 页眉页脚风格
\pagestyle{plain}
}

%% 有时会出现\headheight too small的warning
\setlength{\headheight}{15pt}

%% 清空当前页眉页脚的默认设置
%\fancyhf{}
%%%% fancyhdr设置结束


\makeatletter % 对 beamer 要重新设置
\@ifclassloaded{beamer}{

}{
%%%% 设置listings宏包用来粘贴源代码
%% 方便粘贴源代码，部分代码高亮功能
\usepackage{listings}

%% 设置listings宏包的一些全局样式
%% 参考http://hi.baidu.com/shawpinlee/blog/item/9ec431cbae28e41cbe09e6e4.html
\lstset{
showstringspaces=false,              %% 设定是否显示代码之间的空格符号
numbers=left,                        %% 在左边显示行号
numberstyle=\tiny,                   %% 设定行号字体的大小
basicstyle=\footnotesize,                    %% 设定字体大小\tiny, \small, \Large等等
keywordstyle=\color{blue!70}, commentstyle=\color{red!50!green!50!blue!50},
                                     %% 关键字高亮
frame=shadowbox,                     %% 给代码加框
rulesepcolor=\color{red!20!green!20!blue!20},
escapechar=`,                        %% 中文逃逸字符，用于中英混排
xleftmargin=2em,xrightmargin=2em, aboveskip=1em,
breaklines,                          %% 这条命令可以让LaTeX自动将长的代码行换行排版
extendedchars=false                  %% 这一条命令可以解决代码跨页时，章节标题，页眉等汉字不显示的问题
}}
\makeatother
%%%% listings宏包设置结束


%%%% 附录设置
\makeatletter % 对 beamer 要重新设置
\@ifclassloaded{beamer}{

}{
\usepackage[title,titletoc,header]{appendix}
}
\makeatother
%%%% 附录设置结束


%%%% 日常宏包设置结束
%%%%%%%%------------------------------------------------------------------------


%%%%%%%%------------------------------------------------------------------------
%%%% 英文字体设置结束
%% 这里可以加入自己的英文字体设置
%%%%%%%%------------------------------------------------------------------------

%%%%%%%%------------------------------------------------------------------------
%%%% 设置常用字体字号，与MS Word相对应

%% 一号, 1.4倍行距
\newcommand{\yihao}{\fontsize{26pt}{36pt}\selectfont}
%% 二号, 1.25倍行距
\newcommand{\erhao}{\fontsize{22pt}{28pt}\selectfont}
%% 小二, 单倍行距
\newcommand{\xiaoer}{\fontsize{18pt}{18pt}\selectfont}
%% 三号, 1.5倍行距
\newcommand{\sanhao}{\fontsize{16pt}{24pt}\selectfont}
%% 小三, 1.5倍行距
\newcommand{\xiaosan}{\fontsize{15pt}{22pt}\selectfont}
%% 四号, 1.5倍行距
\newcommand{\sihao}{\fontsize{14pt}{21pt}\selectfont}
%% 半四, 1.5倍行距
\newcommand{\bansi}{\fontsize{13pt}{19.5pt}\selectfont}
%% 小四, 1.5倍行距
\newcommand{\xiaosi}{\fontsize{12pt}{18pt}\selectfont}
%% 大五, 单倍行距
\newcommand{\dawu}{\fontsize{11pt}{11pt}\selectfont}
%% 五号, 单倍行距
\newcommand{\wuhao}{\fontsize{10.5pt}{10.5pt}\selectfont}
%%%%%%%%------------------------------------------------------------------------


%% 设定段间距
\setlength{\parskip}{0.5\baselineskip}

%% 设定行距
\linespread{1}


%% 设定正文字体大小
% \renewcommand{\normalsize}{\sihao}

%制作水印
\RequirePackage{draftcopy}
\draftcopyName{XTUMESH}{100}
\draftcopySetGrey{0.90}
\draftcopyPageTransform{40 rotate}
\draftcopyPageX{350}
\draftcopyPageY{80}

%%%% 个性设置结束
%%%%%%%%------------------------------------------------------------------------


%%%%%%%%------------------------------------------------------------------------
%%%% bibtex设置

%% 设定参考文献显示风格
% 下面是几种常见的样式
% * plain: 按字母的顺序排列，比较次序为作者、年度和标题
% * unsrt: 样式同plain，只是按照引用的先后排序
% * alpha: 用作者名首字母+年份后两位作标号，以字母顺序排序
% * abbrv: 类似plain，将月份全拼改为缩写，更显紧凑
% * apalike: 美国心理学学会期刊样式, 引用样式 [Tailper and Zang, 2006]

\makeatletter
\@ifclassloaded{beamer}{
\bibliographystyle{apalike}
}{
\bibliographystyle{unsrt}
}
\makeatother


%%%% bibtex设置结束
%%%%%%%%------------------------------------------------------------------------

%%%%%%%%------------------------------------------------------------------------
%%%% xeCJK相关宏包

\usepackage{xltxtra,fontspec,xunicode}
\usepackage[slantfont, boldfont]{xeCJK} 

%% 针对中文进行断行
\XeTeXlinebreaklocale "zh"             

%% 给予TeX断行一定自由度
\XeTeXlinebreakskip = 0pt plus 1pt minus 0.1pt

%%%% xeCJK设置结束                                       
%%%%%%%%------------------------------------------------------------------------

%%%%%%%%------------------------------------------------------------------------
%%%% xeCJK字体设置

%% 设置中文标点样式，支持quanjiao、banjiao、kaiming等多种方式
\punctstyle{kaiming}                                        
                                                     
%% 设置缺省中文字体
\setCJKmainfont[BoldFont={Adobe Heiti Std}, ItalicFont={Adobe Kaiti Std}]{Adobe Song Std}   
%% 设置中文无衬线字体
\setCJKsansfont[BoldFont={Adobe Heiti Std}]{Adobe Kaiti Std}  
%% 设置等宽字体
\setCJKmonofont{Adobe Heiti Std}                            

%% 英文衬线字体
\setmainfont{DejaVu Serif}                                  
%% 英文等宽字体
\setmonofont{DejaVu Sans Mono}                              
%% 英文无衬线字体
\setsansfont{DejaVu Sans}                                   

%% 定义新字体
\setCJKfamilyfont{song}{Adobe Song Std}                     
\setCJKfamilyfont{kai}{Adobe Kaiti Std}
\setCJKfamilyfont{hei}{Adobe Heiti Std}
\setCJKfamilyfont{fangsong}{Adobe Fangsong Std}
\setCJKfamilyfont{lisu}{LiSu}
\setCJKfamilyfont{youyuan}{YouYuan}

%% 自定义宋体
\newcommand{\song}{\CJKfamily{song}}                       
%% 自定义楷体
\newcommand{\kai}{\CJKfamily{kai}}                         
%% 自定义黑体
\newcommand{\hei}{\CJKfamily{hei}}                         
%% 自定义仿宋体
\newcommand{\fangsong}{\CJKfamily{fangsong}}               
%% 自定义隶书
\newcommand{\lisu}{\CJKfamily{lisu}}                       
%% 自定义幼圆
\newcommand{\youyuan}{\CJKfamily{youyuan}}                 

%%%% xeCJK字体设置结束
%%%%%%%%------------------------------------------------------------------------

%%%%%%%%------------------------------------------------------------------------
%%%% 一些关于中文文档的重定义
\newcommand{\chntoday}{\number\year\,年\,\number\month\,月\,\number\day\,日}
%% 数学公式定理的重定义

%% 中文破折号，据说来自清华模板
\newcommand{\pozhehao}{\kern0.3ex\rule[0.8ex]{2em}{0.1ex}\kern0.3ex}

\newtheorem{example}{例}                                   
\newtheorem{theorem}{定理}[section]                         
\newtheorem{definition}{定义}
\newtheorem{axiom}{公理}
\newtheorem{property}{性质}
\newtheorem{proposition}{命题}
\newtheorem{lemma}{引理}
\newtheorem{corollary}{推论}
\newtheorem{remark}{注解}
\newtheorem{condition}{条件}
\newtheorem{conclusion}{结论}
\newtheorem{assumption}{假设}

\makeatletter %
\@ifclassloaded{beamer}{

}{
%% 章节等名称重定义
\renewcommand{\contentsname}{目录}     
\renewcommand{\indexname}{索引}
\renewcommand{\listfigurename}{插图目录}
\renewcommand{\listtablename}{表格目录}
\renewcommand{\appendixname}{附录}
\renewcommand{\appendixpagename}{附录}
\renewcommand{\appendixtocname}{附录}
%% 设置chapter、section与subsection的格式
\titleformat{\chapter}{\centering\huge}{第\thechapter{}章}{1em}{\textbf}
\titleformat{\section}{\centering\sihao}{\thesection}{1em}{\textbf}
\titleformat{\subsection}{\xiaosi}{\thesubsection}{1em}{\textbf}
\titleformat{\subsubsection}{\xiaosi}{\thesubsubsection}{1em}{\textbf}

\@ifclassloaded{book}{

}{
\renewcommand{\abstractname}{摘要}
}
}
\makeatother

\renewcommand{\figurename}{图}
\renewcommand{\tablename}{表}

\makeatletter
\@ifclassloaded{book}{
\renewcommand{\bibname}{参考文献}
}{
\renewcommand{\refname}{参考文献} 
}
\makeatother

\floatname{algorithm}{算法}
\renewcommand{\algorithmicrequire}{\textbf{输入:}}
\renewcommand{\algorithmicensure}{\textbf{输出:}}

%%%% 中文重定义结束
%%%%%%%%------------------------------------------------------------------------


\title{有限差分法}
%\author{}
\date{\chntoday}

\begin{document}
%得到相应的线性方程组
%$$
%\textbf{Au}=\textbf{f}
%$$
%矩阵 $A$ 仍然是对称的，但只有半定的。$\textbf{A}$ 的核: $\textbf{Au}=0$ 当且仅当 $\textbf{u}=c$.这需要兼容条件($5$)的离散格式：
%\begin{equation}
%\sum_{i=1}^N f_i=0
%\end{equation}

\section{误差估计}
为了分析误差，我们需要把这个问题变成一个正常的空间。有限维线性空间 $V_h$ 的范数是最大范数：对于$v\in V_h$,
$$
\left \| v\right \|_{\infty,\Omega_h}=\max_{1\le i\le n+1,1\le j\le m+1}\lbrace |v_{i,j}|\rbrace
$$
下标 $h$ 表示这个范数依赖于三角剖分，因为对于不同的 $h$，我们有不同的 $v_{i,j}$.注意，这是 $R^N$ 的 $l^\infty$范数。

\textbf{定理3.1(离散最大值原理)}
令 $v\in V_h$ 满足
$$
\Delta _h v\ge0
$$
那么
$$
\max_{\Omega_h}v\le\max_{\Gamma_h}v
$$
等式成立当且仅当 $v$ 是常数。

证明：假设 $\max_{\Omega_h}v >\max_{\Gamma_h}v$,那么我们可以取一个内部节点 $x_0$,并且在该点达到最大值。设$x_1,x_2,x_3$ 和 $x_4$ 是 $x_0$ 的四个周围点。那么
$$
4v(x_0)=\sum_{i=1}^4 v(x_i)-h^2\Delta _h v(x_0)\le\sum_{i=1}^4 v(x_i)\le 4v(x_0)
$$
等式始终成立，并且由 $4v(x_0)\le\sum_{i=1}^4 v(x_i)$ 可知，$v(x_0)=v(x_i),i=1,2,3,4$,如果一点的值小于 $v(x_0)$,则必有另一点的值大于 $v(x_0)$,与 $x_0$ 处取得最大值矛盾，因此 $v$ 在 $x_0$ 的临近点处也取得最大值。再通过对 $x_0$ 的临近点取临近点，运用上面的方法，最终可以断定 $v$是常数，这与假设 $\max_{\Omega_h}v >\max_{\Gamma_h}v$ 矛盾。

\textbf{定理3.2}
令 $u_h$ 是下面方程的解
\begin{equation}
-\Delta_h u_h=f_I ~ at ~\Omega_h\setminus\Gamma_h,~ u_h=g_I ~ at ~ \Gamma_h
\end{equation}
那么
\begin{equation}
\left \| u_h\right \|_{\infty,\Omega_h}\le\frac{1}{8}\left \| f_I\right \|_{\infty,\Omega_h\setminus\Gamma_h}+\left \| g_I\right \|_{\Gamma_h,\infty}
\end{equation}

证明：我们引入比较函数，
$$
\phi=\frac{1}{4}\left[(x-\frac{1}{2})^2+(y-\frac{1}{2})^2\right]
$$
满足 $\Delta_h\phi_I=1$ 在 $\Omega_h\setminus\Gamma_h$,并且 $0\le\phi\le 1/8$.令 $M=\parallel f_I\parallel _{\infty,\Omega_h\setminus\Gamma_h}$,那么
$$
\Delta_h(u_h+M\phi_I)=\Delta_h u_h+M\ge 0
$$
所以，
$$
\max_{\Omega_h}u_h\le\max_{\Omega_h}(u_h+M\phi_I)\le\max_{\Gamma_h}(u_h+M\phi_I)\le\max_{\Gamma_h}g_I+\frac{1}{8}M
$$
所以，$u_h$ 由($2$)限定。

\textbf{推论3.3}
设 $u$ 是 $Dirichlet$ 问题的真解，$u_h$ 是离散问题($1$)的解。那么
$$
\left \| u_I-u_h\right \|_{\infty,\Omega_h}\le\frac{1}{8}\left \|\Delta_h u_I-(\Delta u)_I\right \|_{\infty,\Omega_h\setminus \Gamma_h}
$$
下面研究误差 $\left \|\Delta_h u_I-(\Delta u)_I\right \|_{h,\infty}$ .

\textbf{引理3.4}
如果 $u\in C^4(\Omega)$,那么
$$
\left \|\Delta_h u_I-(\Delta u)_I\right \|_{\infty,\Omega_h \setminus \Gamma_h}\le\frac{h^2}{6}max\left \{\left \|\frac{\partial ^4u}{\partial x^4}\right \|_{\infty,\Omega},\left \|\frac{\partial ^4u}{\partial y^4}\right \|_{\infty}\right \}
$$
证明：
$$
u(x_{i+1},y_j)=u(x_i,y_j)+h_x\frac{\partial u(x_i,y_j)}{\partial x}+\frac{h_x^2}{2!}\frac{\partial ^2 u(x_i,y_j)}{\partial x^2}+\frac{h_x^3}{3!}\frac{\partial ^3 u(x_i,y_j)}{\partial x^3}+\frac{h_x^4}{4!}\frac{\partial ^4 u(x_i,y_j)}{\partial x^4}
$$

$$
u(x_{i-1},y_j)=u(x_i,y_j)-h_x\frac{\partial u(x_i,y_j)}{\partial x}+\frac{h_x^2}{2!}\frac{\partial ^2 u(x_i,y_j)}{\partial x^2}-\frac{h_x^3}{3!}\frac{\partial ^3 u(x_i,y_j)}{\partial x^3}+\frac{h_x^4}{4!}\frac{\partial ^4 u(x_i,y_j)}{\partial x^4}
$$

因此
$$
\frac{u(x_{i+1},y_j)-2u(x_i,y_j)+u(x_{i-1},y_j)}{h_x^2}=\frac{\partial ^2 u(x_i,y_j)}{\partial x^2}+\frac{h_x^2}{12}\frac{\partial ^4 u(x_i,y_j)}{\partial x^4}
$$
同理可得
$$
\frac{u(x_i,y_{j+1})-2u(x_i,y_j)+u(x_i,y_{j-1})}{h_x^2}=\frac{\partial ^2 u(x_i,y_j)}{\partial y^2}+\frac{h_y^2}{12}\frac{\partial ^4 u(x_i,y_j)}{\partial y^4}
$$
因此有
$$\frac{h_x^2}{12}\frac{\partial ^4 u(x_i,y_j)}{\partial x^4}+\frac{h_y^2}{12}\frac{\partial ^4 u(x_i,y_j)}{\partial y^4}\le\frac{h^2}{12}(\frac{\partial ^4 u(x_i,y_j)}{\partial x^4}+\frac{\partial ^4 u(x_i,y_j)}{\partial y^4})
$$





\textbf{定理3.5}
设 $u$ 是 $Dirichlet$ 问题的真解，以及离散问题($1$)的解。如果 $u\in C^4(\Omega)$,那么
$$
\left \| u_I-u_h\right \|_{\infty,\Omega_h}\le Ch^2
$$
常数 $C$
$$
C=\frac{1}{48}max\left \{\left \|\frac{\partial ^4u}{\partial x^4}\right \|_{\infty,\Omega},\left \|\frac{\partial ^4u}{\partial y^4}\right \|_{\infty}\right \}
$$
\textbf{定理3.5}可由\textbf{推论3.3}和\textbf{引理3.4}得到。

\section{单元中心有限差分法}
在一些实际应用中，比如计算流体动力学 $(CFD)$，泊松方程是在稍微不同的网格上求解的（以单元为中心的均匀网格），如下图。在这一节中，我们
考虑单元处泊松方程的有限差分。

\begin{figure}[H]
\centering
\includegraphics[scale=0.5]{./figures/1.png}
\caption{}
\end{figure}

在内部节点，可以使用标准模板 $(4,−1,−1,−1,−1)$，但靠近边界的节点不能采用上面的模板。内节点之间
的轴向距离仍为 $h$，但靠近边界的节点（单元接触边界的中心）与边界之间的距离为 $h/2$.可以证明，对于 $Neumann$ 边界条件，靠近边界节点的模板是 $(3,−1,−1,−1)$ 和角单元节点是 $(2,−1,−1)$.

处理单元中心差分的 $Dirichlet$ 边界条件时，我们仍然可以引入鬼网格点，并使用标准的 $(4,−1,−1,−1)$ 模
板用于靠近边界的节点。边界上并没有网格点，鬼点处的值需要消除，即需要 $(u_{0,j}+u_{1,j})/2=g_{0,j}:=g_{1/2,j}$.
$$
\frac{5u_{1,j}-u_{2,j}-u_{1,j-1}-u_{1,j+1}}{h^2}=f_{1,j}+\frac{2g_{1/2,j}}{h^2}
$$
对于靠近边界节点，模板是 $(5,-1,-1,-1,-2)$；对于角单元节点，模板是 $(6,-1,-1,-2,-2)$.

但是这种处理是低阶的，为了获得更好的截断误差，可以采用下面的方法，用 $u_{1/2,j}$,$u_{1,j}$,$u_{2,j}$ 处理 $u_{0,j}$.即用 $u_{1/2,j}$,$u_{1,j}$,$u_{2,j}$ 构造一个二次拉格朗日插值多项式，令 $X=X(x,y)$ 表示点的坐标，因此有

$$
l_{1/2,j}(X)=\frac{(X-X_{1,j})(X-X_{2,j})}{(X_{1/2,j}-X_{1,j})(X_{1/2,j}-X_{2,j})}
$$

$$
l_{1,j}(X)=\frac{(X-X_{1/2,j})(X-X_{2,j})}{(X_{1,j}-X_{1/2,j})(X_{1,j}-X_{2,j})}
$$

$$
l_{2,j}(X)=\frac{(X-X_{1/2,j})(X-X_{1,j})}{(X_{2,j}-X_{1/2,j})(X_{2,j}-X_{1,j})}
$$


$$
L_2(X)=u_{1/2,j} l_{1/2,j}(X)+u_{1,j} l_{1,j}(X)+u_{2,j} l_{2,j}(X)
$$

把 $u_{0,j}$ 的坐标带入上式，得到
$$
u_{0,j}=-2u_{1,j}+\frac{1}{3}u_{2,j}+\frac{8}{3}u_{1/2,j}
$$
修改后的方案
$$
\frac{ 6u_{1,j}-\frac{4}{3}u_{2,j}-u_{1,j-1}-u_{1,j+1} }{h^2}=f_{1,j}+\frac{ \frac{8}{3} g_{1/2,j} }{h^2}
$$

我们用 $(6,−\frac{4}{3},−1,−1,−\frac{8}{3})$ 表示靠边界节点的模板。由于截断误差的改善，上面的方法可以提高收敛速度，但是破坏了矩阵的对称性。

对于泊松方程，有一种方法可以保持截断误差和对称性。为简单起见，我们考虑齐次 $Dirichlet$边界条件，即 $u|_{\partial\Omega}=0$,沿边界的切向导数 $\partial ^2_t u=0$.假设方程 $-\Delta u= f$ 也在边界上成立。 注意在边界上，
$\Delta$ 运算符可以写成 $\partial ^2_t+\partial ^2_n$,然后我们在 $\partial\Omega$ 上得到 $\partial ^2_n u=\pm f$.

然后，我们可以用 $u_1,u_{1/2}=0$ 和 $\partial ^2_n=f$ 该函数，并得到鬼点处的一个方程
$$
u_{1,j}+u_{0,j}=\frac{h^2}{4}f_{1/2,j}
$$
并将边界模板修改为
$$
\frac{5u_{1,j}-u_{2,j}-u_{1,j-1}-u_{1,j+1}}{h^2}=f_{1,j}+\frac{1}{4}f_{1/2,j}
$$













%坎坎坷坷扩


%\cite{tam19912d}
%\bibliography{../ref}
\end{document}
